% changelog.tex
%
% The documentation in this file is part of Pyxplot
% <http://www.pyxplot.org.uk>
%
% Copyright (C) 2006-2026 Dominic Ford <coders@pyxplot.org.uk>
%               2008-2013 Ross Church
%
% $Id$
%
% Pyxplot is free software; you can redistribute it and/or modify it under the
% terms of the GNU General Public License as published by the Free Software
% Foundation; either version 2 of the License, or (at your option) any later
% version.
%
% You should have received a copy of the GNU General Public License along with
% Pyxplot; if not, write to the Free Software Foundation, Inc., 51 Franklin
% Street, Fifth Floor, Boston, MA  02110-1301, USA

% ----------------------------------------------------------------------------

% LaTeX source for the Pyxplot Users' Guide

\chapter{ChangeLog}
\index{ChangeLog}

\subsection*{2026 Nov 19: Pyxplot 0.10.0}

Version 0.10.0 is a major refactoring of Pyxplot's source code, which has been reformatted, documented throughout, and made const-correct. It also fixes a number of bugs, including:

\begin{itemize}
\item Dates displayed and entered in the Hebrew and Islamic calendars were incorrect.
\item Setting a preferred unit in a configuration file made Pyxplot run out of memory on startup.
\item The commands {\tt unset unit preferred} and {\tt unset unit of} did not work.
\item The help text for several functions gave incorrect names.
\item Several memory errors, including a use-after-free in the {\tt histogram} command.
\item Crashes in the {\tt conjugate()} function, in pie charts of unlabelled data supplied inline, and when an arrow or label was positioned with units incompatible with its axis.
\item Macro substitution ({\tt @variable}) removed quotation marks from commands.
\item {\tt texify()} inserted an empty {\tt \textbackslash mathrm\{\}} after function calls.
\item Dictionary items without values, e.g.\ {\tt \{'a':1, 'b'\}}, and ternary operators without a colon, e.g.\ {\tt 1 ? 2}, gave misleading error messages instead of syntax errors.
\item Warnings about the ranges of logarithmic axes were given twice.
\item {\tt acosh()} returned zero for all real arguments.
\item {\tt besseli()} and {\tt besselk()} returned exponentially scaled values.
\item Type objects all compared as equal, e.g.\ {\tt types.number == types.string}.
\item {\tt time.intervalStr()} gave negative intervals when a format was given.
\item {\tt time.fromCalendar()} discarded fractions of a second, and dates lost microseconds of precision when converted back to calendar dates.
\item {\tt time.now()} had a resolution of only one second.
\item After {\tt os.chdir()} was given a relative path, {\tt os.getcwd()} returned it unchanged, and Pyxplot returned to the wrong directory after plotting.
\item {\tt phy.Bvmax()} used an imprecise value of Wien's constant.
\item {\tt texify()} produced unbalanced \latexdcf\ for some functions, including {\tt range()} and the hypergeometric functions.
\item The help text of {\tt sinc()} described the normalised sinc function.
\item {\tt sgn()} did not accept quantities with physical units.
\item The {\tt toMinute()} method of dates returned the day of the month, and {\tt toDayWeekNum()} returned negative numbers for dates before 1970.
\item The {\tt componentsCMYK()} method of colors returned the black component in place of the yellow one.
\item The {\tt eigenvectors()} method of matrices returned the rows of the matrix of eigenvectors rather than its columns, so did not return the eigenvectors. {\tt eigenvalues()} and {\tt eigenvectors()} now both list their results in order of ascending eigenvalue.
\item The {\tt setPos()} method of files rejected the positions returned by {\tt getPos()}, which have units of bytes.
\item The index blocks written by the {\tt tabulate} command were labelled from 1, though the {\tt index} modifier counts them from 0.
\item The {\tt show} command described minor ticks as major ticks.
\end{itemize}

\noindent The unit {\tt planck\_impedence} has been renamed {\tt planck\_impedance}; the old spelling is still accepted. The documentation has been checked against the source code: commands, settings and configuration-file options that were missing have been documented, and errors in the names, defaults and syntax given for others have been corrected.

\noindent Pyxplot now exits with status 1 if any error occurs while running a script or piped input, so that failures can be detected by the shell. The lengths, indices and slices of strings, and the positions returned by the {\tt find()} and {\tt findAll()} methods, now count characters rather than bytes, so that multi-byte UTF-8 characters such as {\tt \'e} count as one character, and the {\tt isalpha()} and {\tt isalnum()} methods recognise letters in any alphabet. An automated test suite has been added, which may be run with {\tt make test}. When compiled with OpenMP, the {\tt plot}, {\tt tabulate} and {\tt fit} commands evaluate functions on all of the computer's cores at once, when this gives the same results as evaluating them one at a time: e.g.\ not when they use random numbers or call subroutines. The number of cores used can be limited by setting the environment variable {\tt OMP\_NUM\_THREADS}. Fits to more than 4096 data points sum their residuals in blocks, so their results may differ in the last digits from earlier versions, but do not depend on the number of cores used.

\subsection*{2012 Sep 19: Pyxplot 0.9.2}

Version 0.9.2 corrects a large number of minor bugs.

\subsection*{2012 Aug 29: Pyxplot 0.9.1}

Version 0.9.1 is a minor update with new support for running Pyxplot on Raspberry Pi. It fixes SIGBUS errors in Pyxplot's math engine when run on armhf architectures.

\subsection*{2012 Aug 1: Pyxplot 0.9.0}

Version 0.9 is a major update. Many new data types have been introduced, each
of which has methods which can be called in an object-orientated fashion. These
include:

\begin{itemize}
\item {\bf Colors}, which can be stored in variables for subsequent use in vector graphics commands. The addition and subtraction operators act on colors to allow color mixing.
\item {\bf Dates}, which can be imported from calendar dates, unix times or Julian dates. Dates can be subtracted to give time intervals.
\item {\bf Lists} and {\bf dictionaries}, which can be iterated over, or used to feed calculated data into the {\tt plot} and {\tt tabulate} commands.
\item {\bf Vectors} and {\bf matrices}, which allow matrix algebra. These types interface cleanly with Pyxplot's vector-graphics commands, allowing positions to be specified as vector expressions.
\item {\bf File handles}, which allow Pyxplot to read data from files, or write data or logs to files.
\item {\bf Modules} and {\bf classes}, which allow object-orientated programming.
\end{itemize}

In addition, Pyxplot's range of operators has been extended to include most of those in the C programming language, allowing expressions such as

\vspace{3mm}
\input{fragments/tex/fs_operators.tex}
\vspace{3mm}

\noindent to be written.

\subsubsection*{Incompatibilities with Pyxplot 0.8}

The extensions to Pyxplot in version 0.9 mean that some minor changes to syntax have been necessary. These include:

\begin{itemize}
\item Some functions and variables have been renamed. Variables whose names used to begin {\tt phy\_} now live in a module called {\tt phy}. They may be accessed as, for example, {\tt phy.c}. Similarly, random number generating functions now live in a module called {\tt random}; statistics functions in a module called {\tt stats}; time-handling functions in {\tt time}; operating system functions in {\tt os}; and astronomy functions in {\tt ast}. The contents of these modules can be listed by typing, for example, {\tt print phy}.
\item Custom colors, which used to be specified using syntax such as {\tt rgb0.2:0.3:0.4}, should now be specified using the {\tt rgb(r,g,b)} functions, as, for example, {\tt rgb(0.2,0.3,0.4)}. Custom colors can now be stored in variables for later use (see Section~\ref{sec:colorObjects}).
\item The range of escape characters which can be used in strings has been increased, so that, for example, {\tt $\backslash$n} is a newline and {\tt $\backslash$t} a tab. As in python, prepending the string with the character {\tt r} disables all escape character expansion. As backslashes are common characters in latex command strings, the easiest approach is to always prepend latex strings with an {\tt r}. As in python, triple quotes, e.g.\ {\tt r"""2 $\backslash$times 3"""} can be used where required (see Section~\ref{sec:latex_incompatibility}).
\item In the {\tt foreach} command, square brackets should be used to delimit lists of items to iterate over. The Pyxplot 0.8 syntax {\tt foreach i in (1,2,3)} should now be written {\tt foreach i in [1,2,3]} (see Section~\ref{sec:foreach}).
\end{itemize}

\subsection*{2011 Jan 7: Pyxplot 0.8.4}

\subsubsection*{Summary:}

This is a minor bugfix release.

\subsubsection*{Details:}

\begin{itemize}
\item Two-dimensional parametric grid plotting implemented.
\item Bugfix to the dots plot style; filled triangles replaces with filled circles.
\item Bugfix to linewidths used when drawing line icons on graph legends.
\item Bugfix to Makefile to ensure libraries link correctly under Red Hat and SUSE.
\item Code cleanup to ensure correct compilation with {\tt -O2} optimisation.
\end{itemize}

\subsection*{2010 Sep 15: Pyxplot 0.8.3}

\subsubsection*{Summary:}

This is a minor bugfix release.

\subsubsection*{Details:}

\begin{itemize}
\item @ macro expansion operator implemented.
\item assert command implemented.
\item for command behaviour changed such that {\tt for i=1 to 10} includes a final iteration with {\tt i}=10.
\item Point types rearranged into a more logical order.
\item Improved support for newer Windows bitmap images.
\item Bugfix to the {\tt set unit preferred} command.
\item Binary not operator bugfixed.
\item Bugfix to handling of comma-separated horizontal datafiles.
\item Mathematical function {\tt finite()} added.
\end{itemize}

\subsection*{2010 Aug 4: Pyxplot 0.8.2}

\subsubsection*{Summary:}

This release introduces three-dimensional plotting, as well as the ability to plot two-dimensional maps of functions as either color maps, contour plots, or as three-dimensional surfaces. A large number of bugs have also been fixed.

\subsubsection*{Details:}

\begin{itemize}
\item 3D plotting implemented.
\item New plot styles colormap, contourmap and surface added.
\item Interpolation of 2D datagrids and bitmap images implemented.
\item Stepwise interpolation mode added.
\item Dependency on libkpathsea relaxed to make installation under MacOS easier; linking to the library is still strongly recommended on systems where it is readily available.
\item Mathematical functions {\tt frac-tal\_\-julia()}, {\tt frac\-tal\_\-man\-del\-brot()} and {\tt prime()} added.
\item Many bug fixes, especially to the ticking of axes.
\end{itemize}

\subsection*{2010 Jun 1: Pyxplot 0.8.1}

\subsubsection*{Summary:}

This release has no major new features, but fixes several significant bugs in version 0.8.0.

\subsubsection*{Details:}

\begin{itemize}
\item Mathematical functions {\tt time\_\-from\-unix()}, {\tt time\_\-unix()}, {\tt zernike()} and {\tt zernikeR()} added.
\item Bug fix to the ticking of linked axes.
\item Bug fix to the ticking of axes with blank axis tick labels.
\item Makefile and configure script improved for portability.
\end{itemize}

\subsection*{2010 May 19: Pyxplot 0.8.0}

\subsubsection*{Summary:}

This release is a major update, for which Pyxplot's original python code has
been completely rewritten in C with the addition of many new features. Because
of the scale of this update, there is some minor syntax incompatibility with
previous versions where features have undergone particularly heavy change. The
most apparent change is the increase in speed and efficiency resultant from the
use of a compiled language: especially when handling large \datafile s, Pyxplot
0.8.0 can run more than an order-of-magnitude faster than previous versions.

\subsubsection*{Details:}

\begin{itemize}
\item The handling of large \datafile s has been streamlined to require around an order-of-magnitude less time and memory.
\item Pyxplot's mathematical environment has been extended to handle complex numbers and quantities with physical units.
\item The range of mathematical functions built into Pyxplot has been massively extended.
\item The {\tt solve} command has been added to allow the solution of systems of equations.
\item The {\tt maximize} and {\tt minimize} commands have been added to allow searches for local extrema of functions.
\item An {\tt fft} command has been added for performing Fourier transforms on data.
\item New plot styles -- {\tt filledregion} and {\tt yerrorshaded} -- have been added for plotting filled error regions.
\item The configuration of linked axes has been entirely redesigned.
\item Parametric function plotting has been implemented.
\item Colours can now be specified by RGB, HSB or CMYK components, as well as by name.
\item Several commands, e.g. {\tt box}, {\tt circle}, {\tt ellipse}, etc., have been added to allow vector graphics to be produced in Pyxplot's multiplot environment.
\item The {\tt jpeg} command has been generalised to allow the incorporation of not only {\tt jpeg} images, but also {\tt bmp}, {\tt gif} and {\tt png} images, onto multiplot canvases. The command has been renamed {\tt image} in recognition of its wider applicability. Image transparency is now supported in {\tt gif} and {\tt png} images.
\item The {\tt spline} command, now renamed the {\tt interpolate} command, has been extended up provide many types of interpolation between datapoints.
\item A wide range of conditional and flow control structures have been added to Pyxplot's command language -- these are the {\tt do}, {\tt for}, {\tt foreach}, {\tt if} and {\tt while} commands and the {\tt cond\-ition\-alS} and {\tt con\-dition\-alN} mathematical functions.
\item Input filters have been introduced as a mechanism by which datafiles in arbitrary formats can be read.
\item Pyxplot's command-line interface now supports tab completion.
\item The {\tt show} command has been reworked to produce pastable output.
\item Many minor bugs have been fixed.
\end{itemize}

\subsection*{2009 Nov 17: Pyxplot 0.7.1}

\subsubsection*{Summary:}

This release has no major new features, but fixes several serious bugs in version 0.7.0.

\subsubsection*{Details:}

\begin{itemize}
\item The {\tt exec} command did not work in Pyxplot 0.7.0; this issue has been resolved.
\item The {\tt xyerrorrange} plot style did not work in Pyxplot 0.7.0; this issue has been resolved.
\item Pyxplot 0.7.0 produces large numbers of python deprecation error messages when run under python 2.6; the code has been updated to remove references to deprecated python functions.
\end{itemize}

\subsubsection*{Details -- Change of System Requirements:}

\begin{itemize}
\item In order to fix some of the bugs listed above, it has been necessary to
fix bugs in the PyX graphics library as well as those in Pyxplot. As a result,
and to ensure that these bugfixes reach users as quickly as possible, we have
opted to ship our own modified version of PyX 0.10, called dcfPyX with Pyxplot.
\end{itemize}

\subsection*{2008 Oct 14: Pyxplot 0.7.0}

\subsubsection*{Summary:}

Third Pyxplot beta-release. The code has undergone significant streamlining,
and now runs approximately twice as fast as version 0.6.3 when handling large
datafiles. Memory usage has also been radically reduced. Two new data
processing commands have been introduced. The {\tt tabulate} command can be
used to produce textual datafiles, allowing the user to read data in from
files, apply some analysis, and then write the processed data back to file. The
{\tt histogram} command can be used to estimate the frequency densities of sets
of data points, either by binning them into a bar chart, or by fitting a
functional form to their frequency density.

\subsubsection*{Details -- New and Extended Commands:}

\begin{itemize}
\item {\tt tabulate}
\item {\tt histogram}
\item {\tt set label} and {\tt text} commands extended to allow a color to be
specified.
\end{itemize}

\subsubsection*{Details -- API changes}

\begin{itemize}
\item {\tt diff\_dx()} and {\tt int\_dx()} functions -- the function to be
differentiated or integrated must now be placed in quotation marks.
\end{itemize}

\subsubsection*{Details -- Change of System Requirements:}

\begin{itemize}
\item Requirement of PyX version 0.9 has been updated to PyX version 0.10. Note that new versions of the PyX graphics library are not generally backwardly compatible.
\end{itemize}

\subsection*{2007 Feb 26: Pyxplot 0.6.3}

\subsubsection*{Summary:}

Second Pyxplot beta-release. The most significant change is the introduction of
a new command-line parser, with greatly improved handling of complex
expressions and much more meaningful syntax error messages. Multi-platform
compatibility has also been massively improved, and dependencies loosened.  A
small number of new commands have been added; most notable among them are the
{\tt jpeg} and {\tt eps} commands, which embed images in multiplots.

\subsubsection*{Details -- New and Extended Commands:}

\begin{itemize}
\item {\tt jpeg}
\item {\tt eps}
\item {\tt set xtics} and {\tt set mxtics}
\item {\tt text} and {\tt set label} commands extended to allow text rotation.
\item {\tt set log} command extended to allow the use of logarithms with bases other than 10.
\item {\tt set preamble}
\item {\tt set term enlarge | noenlarge}
\item {\tt set term pdf}
\item {\tt set term x11\_persist}
\end{itemize}

\subsubsection*{Details -- Eased System Requirements:}

\begin{itemize}
\item Requirement on Python 2.4 minimum eased to version 2.3 minimum.
\item Requirements on scipy and readline eased; Pyxplot will now work in reduced form when they are absent, though they are still strongly recommended.
\item dvips and Ghostscript are no longer required.
\end{itemize}

\subsubsection*{Details -- Removed Commands:}

Due to a general refinement of Pyxplot's API, some of the less sensible pieces
of syntax from Version~0.5 are no longer supported. The author apologises for
any inconvenience caused.

\begin{itemize}
\item The {\tt delete\_arrow}, {\tt delete\_text}, {\tt move\_text}, {\tt undelete\_arrow} and {\tt undelete\_text} commands have been removed from the Pyxplot API. The {\tt move}, {\tt delete} and {\tt undelete} commands should now be used to act upon all types of multiplot objects.
\item The {\tt set terminal} command no longer accepts the {\tt enhanced} and {\tt noenhanced} modifiers. The {\tt postscript} and {\tt eps} terminals should be used instead.
\item The {\tt select} modifier, used after the {\tt plot}, {\tt replot}, {\tt fit} and {\tt spline} command can now only be used once; to specify multiple {\tt select} criteria, use the {\tt and} logical operator.
\end{itemize}

\subsection*{2006 Sep 09: Pyxplot 0.5.8}

First beta-release.
